\section{Representative functions and localized radii}\label{sec:local54}
A real representative function on $K$ is a matrix coefficient of a finite-dimensional continuous orthogonal representation, or a finite sum of such coefficients. Their algebra contains the constants, is closed under multiplication by tensor products, and is uniformly dense in $C(K,\R)$ by the Peter--Weyl approximation theorem \cite{Folland}. A finite collection of representative functions can be written on one space
\[
 \mathcal V=\bigoplus_{i=1}^a\End(E_i),\qquad
 \Phi(h)=(R_i(h))_{i=1}^a,
\]
where the scalar product is the sum of Frobenius products. Left multiplication defines an orthogonal representation $R$ on $\mathcal V$ and satisfies $\Phi(gh)=R(g)\Phi(h)$. Adjoin a trivial summand when necessary. Let
\[
 P=\int_KR(h)\,d\mu(h),\quad E=(I-P)\mathcal V,\quad
 Z(h)=(I-P)\Phi(h),\quad Q_0=\norm{Z(e)}.
\]
Haar averaging is the orthogonal projection onto fixed vectors. Consequently $Z(gh)=R(g)Z(h)$, the representation on $E$ has no fixed vectors, and $P\Phi(h)=P\Phi(e)$ for $h\in K$.

Choose and fix a coefficient vector $a_p$ for each scalar representative function $p$, and a linear coefficient map $A_G$ for each vector-valued representative function $G$. Their norms need not be intrinsic; fixing any such choices suffices. For an arbitrary joint distribution of $h\in K$ and a finite state $S$, put $Q=\E\norm{\E[Z(h)\mid S]}$ and $\bar p=\int p\,d\mu$. The fixed/moving decomposition gives
\begin{align}
 \E[p(h)\mid S]&=\bar p+\ip{a_p}{\E[Z(h)\mid S]},\label{eq:calibration54}\\
 |\E p-\bar p|&\le\norm{a_p}Q,\qquad
 \norm{\E G-\bar G}\le\norm{A_G}Q.\label{eq:vectorcal54}
\end{align}
Here the norm of $A_G$ is the operator norm from the Euclidean feature space to the indicated output space. If $D:S\to Y$ and $M_Y=\max_{y\in Y}\norm{y}$, then
\begin{equation}\label{eq:decodercal54}
 \norm{\E[p(h)D(S)]-\bar p\E D(S)}\le M_Y\norm{a_p}Q.
\end{equation}
Indeed condition on $S$ in \eqref{eq:calibration54}, multiply by $D(S)$, and apply the triangle inequality. These are moment estimates, not total-variation convergence of a hidden physical distribution.

\begin{lemma}[Finite localization of a radius]\label{lem:radiuslocal54}
Let $f:K\to Y$ be continuous, with $Y$ compact and convex in a finite-dimensional normed space. Suppose $\rad_Y(f(K))>\varepsilon$. There are a vector-valued representative function $G$, a number $\eta>0$, and nonnegative representative functions $p_1,\ldots,p_a$ of Haar mean one such that
\[
 \norm{G-f}_\infty<\eta,\qquad
 r_*:=\min_{y\in Y}\max_i\norm{\bar{p_iG}-y}>\delta:=\varepsilon+\eta.
\]
\end{lemma}
\begin{proof}
Choose $\varepsilon<t<\rad_Y(f(K))$. For each $y\in Y$ some $h\in K$ has $\norm{f(h)-y}>t$. The corresponding open subsets of $Y$ cover $Y$, so a finite subcover yields $h_1,\ldots,h_a$ with
\[
 \min_{y\in Y}\max_i\norm{f(h_i)-y}>t.
\]
The strict inequality follows also at a minimizing $y$, since the maximum is everywhere greater than $t$ and $Y$ is compact.

Fix a sufficiently small positive $\eta$. Approximate the coordinates of $f$ uniformly by representative functions to obtain $\norm{G-f}_\infty<\eta$. For each $i$, choose a nonzero nonnegative continuous bump supported in a neighborhood where $f$ differs from $f(h_i)$ by less than $\eta$. Such a bump exists by normality of the compact Hausdorff space; it has positive Haar integral by full support. Approximate its square root uniformly by a representative function, square the approximation, and normalize its Haar integral. The resulting nonnegative representative function $p_i$ has mean one, and its weighted $f$-mean can be made arbitrarily close to that of the bump. Thus we may ensure
\[
 \norm{\bar{p_iG}-f(h_i)}<3\eta.
\]
The constrained radius changes by at most the maximum perturbation of the finitely many points. Taking $4\eta<t-\varepsilon$ gives $r_*>\varepsilon+\eta$.
\end{proof}

\begin{proposition}[Residual moment from a vector interface]\label{prop:residual54}
Fix the data of Lemma~\ref{lem:radiuslocal54}. Choose one feature space representing all $p_i$ and $p_iG$, and fix their coefficient choices before defining
\[
 L=\max_i\bigl(\norm{A_{p_iG}}+(M_Y+\delta)\norm{a_{p_i}}\bigr),
 \qquad \Delta=r_*-\delta>0.
\]
Suppose a terminal decoder $D(S)\in Y$ obeys
\[
 \norm{\E[D(S)\mid w]-G(h(w))}\le\delta
\]
for every complete deterministic word in a test law. Then $L>0$, $E\ne\{0\}$, $Q_0>0$, and
\begin{equation}\label{eq:residual54}
                         Q\ge\Delta/L.
\end{equation}
\end{proposition}
\begin{proof}
For $p=p_i\ge0$, wordwise correctness and the triangle inequality imply
\[
 \norm{\E[p(h)D(S)]-\E[p(h)G(h)]}\le\delta\E p(h).
\]
This is a conditional mean error, not a pointwise bound on the random decoder. Write $y=\E D(S)$, which belongs to $Y$ by convexity. Equations~\eqref{eq:vectorcal54}--\eqref{eq:decodercal54}, together with $\bar p=1$, give
\[
 \norm{y-\bar{p_iG}}\le\delta+
 \bigl(\norm{A_{p_iG}}+(M_Y+\delta)\norm{a_{p_i}}\bigr)Q.
\]
Maximizing in $i$ yields $r_*\le\delta+LQ$. Hence $LQ\ge\Delta>0$. In particular the moving feature and its initial norm cannot vanish: by equivariance, $Q_0=0$ would imply $Z(h)=0$ for all $h$ and therefore $Q=0$.
\end{proof}

\subsection{Proof of Theorem~\ref{thm:radius54}}
Choose $(s,j,c)$ with component radius greater than $\varepsilon$. By density there is an executable prefix $v$ with $u_v\in c$. Put $b=|v|$ and $f(h)=F_{s,j}(hu_v)$ on $K$. Since $Ku_v=c$, this map has the same constrained radius. Apply Lemma~\ref{lem:radiuslocal54} and Proposition~\ref{prop:residual54}.

The moving feature is nonzero independently of any particular test law. Otherwise all represented functions would be constant on $K$; since each $p_i$ has mean one, this would force $p_i=1$ and $G$ constant. Then all $\bar{p_iG}$ would equal a point at distance less than $\eta$ from $Y$, contradicting $r_*>\delta\ge\eta$.

For a fixed width $k$, apply Lemma~\ref{lem:gap54} to the resulting moving feature. Let $B,d$ be its block parameters and set
\[
 \chi=-\log(1-d)>0,\qquad A=\max(1,LQ_0/\Delta).
\]
Start a test word with $v$; between independent length-$B$ test blocks use fixed return words. The accumulated physical product is $hu_v$. Initially $h=e$, so the conditional feature norm after the prefix equals $Q_0$, regardless of the hidden-state distribution. If $J$ block endpoints have width at most $k$, Lemma~\ref{lem:centroid54} gives
\[
 Q_N\le Q_0(1-d)^J.
\]
The machine is correct on every deterministic word of length $N$. Thus any external law supported on such words inherits the conditional mean bound, with error $\delta=\varepsilon+\eta$ for $G$. Proposition~\ref{prop:residual54} implies
\begin{equation}\label{eq:J54}
                         J\chi\le\log A.
\end{equation}
The law depends on the advertised width profile, not on the machine's private state. When $A=1$, \eqref{eq:J54} excludes even one selected narrow block; no logarithmic loss has occurred.

For clarity we give the spacing count with the full return hypothesis. Discard narrow cuts before $b+B+g$ and after $N-g$. From the remaining cuts greedily choose the first and then the next at distance at least $B+g$. Put an independent block immediately before each selected cut. The initial gap, all intervening gaps, and the final suffix have lengths at least $g$, so \eqref{eq:returns54} fills them with executable returns. At most $b+B+2g-1$ cuts were discarded. Each chosen endpoint accounts for at most $B+g$ eligible cuts. Therefore one may take
\begin{equation}\label{eq:Lk54}
 L_k=\left\lceil b+B+2g-1+(B+g)\frac{\log A}{\chi}\right\rceil.
\end{equation}
When the eligible interval is empty the same estimate follows directly by counting. This proves occupation without restricting any register inside a block or gap. If the entire peak is at most $k$, then $N\le L_k$, which proves divergence directly. Monotonicity in the horizon is not needed when an identity letter is absent.

For the upper bound, use labels $(s,c)\in\mathcal S\times C$, including any labels not reached at a given horizon. Initialize at $(s,K)$ and update by left multiplication
\[
 (s,c)\longmapsto(s,[u_a]c).
\]
Choose a minimizing center of $F_{s,j}(c)$ in $Y_j$ as decoder at $(s,c)$ for query $j$. These are legal decoder values, and the error is at most \eqref{eq:radius54}. All command updates are permutations of one finite state set, independent of clock and horizon. This proves boundary attainment and the equivalence.\qed

\begin{proposition}[A finite return certificate]\label{prop:return54}
Condition~\eqref{eq:returns54} follows if there are finitely many executable identity-product words whose positive lengths have greatest common divisor one.
\end{proposition}
\begin{proof}
Let the lengths be $\ell_1,\ldots,\ell_a$. Their nonnegative sums modulo $\ell_1$ form the subgroup generated by their residues: a subsemigroup of a finite group is a subgroup. The gcd hypothesis makes this the whole residue group. Choose one nonnegative sum $n_r$ for each residue $r$ modulo $\ell_1$, and put $g=\max_r n_r$. Every $n\ge g$ equals its representative $n_r$ plus a nonnegative multiple of $\ell_1$. Concatenating the corresponding return words proves the claim.
\end{proof}
For example, the alphabet $\{R_\theta,R_{-\theta},R_{2\pi/3}\}$, with $\theta/2\pi$ irrational, contains no identity letter but has return lengths two and three. It satisfies \eqref{eq:returns54} with $g=2$. A single irrational rotation letter does not: there is only one word at each horizon, and a horizon-specific constant decoder can store its numerical answer with one command label. Thus synchronization cannot be silently omitted. Nor is a bounded irreversible semigroup enough: for the identity command $I=1$ and contraction command $a=\lambda\in(0,1)$, the mean $\rho\lambda^{\#a(w)}$ has an exact alive/absorbed two-state realization. The group generated with inverses is not bounded. The present theorem is not a compact-semigroup classification.
